\section{Erklärung zum Einsatz von KI‑Systemen}

Gemäß den Richtlinien der Dualen Hochschule Baden‑Württemberg muss in wissenschaftlichen Arbeiten offengelegt werden, ob und wie KI‑basierte Systeme genutzt wurden. In der vorliegenden Arbeit wurde ein Sprachmodell (ChatGPT) zur Unterstützung in folgenden Bereichen verwendet:

% TODO: Offizielle DHBW-Quelle zur KI-Offenlegung konkret nachtragen.

\begin{itemize}[leftmargin=1.2cm]
  \item \textbf{Ideengenerierung}: Bei der Strukturierung der Gliederung und der Entwicklung von Überschriften.
  \item \textbf{Formulierungsunterstützung}: Beim Entwurf einzelner Kapitelabschnitte, um eine klare und präzise Sprache zu gewährleisten.
  \item \textbf{Überprüfung der Logik}: Durch das Erstellen von Listen und Tabellen zur Anforderungsanalyse und Ergebnisdarstellung.
\end{itemize}

Das KI‑System diente ausschließlich als Hilfsmittel. Die inhaltliche Verantwortung, die Auswahl der Quellen, die Interpretation der Ergebnisse und die Endredaktion lagen vollständig beim Autor. Alle Texte wurden kritisch geprüft und nachbearbeitet. Die Nutzung von ChatGPT stellt eine Form der Schreibassistenz dar und begründet keinen Verzicht auf die eigenständige wissenschaftliche Leistung.